\chapter*{Prólogo — 48 horas de blackout}
\addcontentsline{toc}{chapter}{Prólogo — 48 horas de blackout}

A cidade apagou às 18h42 de uma terça-feira quente. Primeiro veio o silêncio estranho dos aparelhos desligando ao mesmo tempo. Depois, um coro de apitos dos nobreaks. A luz piscou, morreu. A tela do celular continuou brilhando por alguns minutos, mas o 4G sumiu em seguida. As mensagens pararam de ser entregues. O aplicativo do banco travou; o \textit{Pix} não ia, o cartão não passava. A portaria do prédio comunicou pelo interfone que os elevadores foram travados por segurança.

Às 19h15, você foi à janela. Uma rua inteira de janelas escuras e lanternas vacilantes. O mercado da esquina levantou meia porta para ventilar. Algumas pessoas corriam com sacolas para dentro de casa. Outras discutiam com o caixa, que não sabia como registrar as compras sem o sistema. Um vizinho tentou usar dinheiro vivo. O caixa não tinha troco.

No primeiro dia, a maior dificuldade foi \textbf{Narrativa}. Ninguém sabia ao certo o que estava acontecendo. O boato mais compartilhado no prédio: “\textit{É ataque cibernético}”. Outro: “\textit{É racionamento de energia}”. Um terceiro: “\textit{É o fim do mundo}”. Sem internet, o condomínio improvisou um quadro de avisos no hall com recados escritos à mão. As pessoas queriam \textit{uma} história que explicasse o caos.

Na manhã seguinte, faltou água no 10º andar: as bombas do prédio dependem de energia. A geladeira começou a descongelar. O segundo golpe foi de \textbf{Identidade}. Sem internet, logins não funcionavam; sem login, muitos serviços não reconheciam você. Um vizinho preso do lado de fora do trabalho não conseguiu passar pela catraca porque o acesso é biométrico na nuvem. Outro não conseguiu “abrir um chamado” para a empresa de luz: o app exigia autenticação de dois fatores\footnote{Dois fatores são ótimos---quando existem os dois fatores.}.

À noite, a tensão subiu. Para vários comércios, o problema se tornou \textbf{Controle}. Sem pagamentos digitais, sem sistema de estoque, sem segurança privada (as câmeras também caíram), alguns fecharam as portas. Outros venderam apenas por anotações em cadernos “para acertar depois”. A padaria limitou a compra de pães por família. Alguém sugeriu um “rateio” de água para o prédio. Um grupo marcou oração no térreo: sem som, sem telão, sem QR Code; só vozes.

\begin{nicbox}
\textbf{Três perguntas NIC para um blackout:}\\[-1ex]
\begin{enumerate}[leftmargin=*,nosep]
  \item \textbf{Narrativa} --- Quem define a história quando os algoritmos emudecem? Como discernir boato de verdade?
  \item \textbf{Identidade} --- Quem somos sem nossos logins, QR Codes e senhas? O que nos confirma quando o sistema não confirma?
  \item \textbf{Controle} --- O que fazemos quando as portas de compra e venda fecham? Como praticar justiça e generosidade quando a escassez bate?
\end{enumerate}
\end{nicbox}

No final das primeiras 24 horas, as famílias que estavam mais tranquilas tinham três traços simples: um plano de comunicação offline, alguma reserva básica (água, gás, dinheiro trocado) e uma comunidade próxima (parentes, vizinhos, igreja). Nada heroico; só \textit{prudência}.

\begin{checklist}[title={Primeiras 24 horas: ações simples}]
\begin{itemize}[leftmargin=*,noitemsep]
  \item Combine \textit{pontos de encontro} e \textit{horários} de checagem presencial com a família.
  \item Desligue tomadas críticas para proteger equipamentos quando a energia voltar.
  \item Separe água potável e identifique recipientes de uso sanitário.
  \item Organize alimentos que não exigem refrigeração; priorize o que pode estragar.
  \item Faça um \textit{mapa de vizinhança}: quem tem necessidades especiais? Quem pode ajudar?
  \item Anote no papel contatos essenciais (família, líderes, serviços).
  \item Defina um \textit{horário fixo} de oração e leitura breve da Palavra.
\end{itemize}
\end{checklist}

No final das 48 horas, a luz voltou em ondas, bairro a bairro. As redes sociais retornaram como uma torrente. As narrativas competiam outra vez, mas agora com gráficos e especialistas. Você respirou aliviado, mas percebeu que algo tinha mudado: o \emph{mundo}---não apenas os aparelhos---ficara frágil.

Este livro começa aqui. Não para alimentar pânico, mas para formar \textbf{prudência, esperança e fidelidade}. Ao longo das próximas páginas, vamos mapear como tecnologia e Apocalipse se cruzam nas \textbf{três frentes NIC}:
\begin{enumerate}[leftmargin=*,nosep]
  \item \textbf{Narrativa} (\textit{verdade}) --- quando a “imagem que fala” cria mundos plausíveis;
  \item \textbf{Identidade} (\textit{quem sou}) --- quando somos reduzidos a perfis e credenciais;
  \item \textbf{Controle} (\textit{o que posso fazer}) --- quando portas de compra e venda são programáveis.
\end{enumerate}

E faremos isso de modo bíblico e prático: com critérios de leitura, cenários realistas, checklists e disciplinas espirituais. Porque, quando a cidade apaga, \textit{a igreja acende}. E quando a nuvem falha, o Cordeiro permanece.

\bigskip
\noindent\textit{Próximo capítulo: Como ler Apocalipse sem delirar.}
